\chapter{The Inventory by Discipline}
\label{ch:04}

An inventory of unsolved problems can be organized in several ways, and the choice of organization is not neutral. Organization by urgency invites a ranking of lives against lives, and organization by technology invites the assumption that the technology is the answer. Organization by discipline has a more modest and more useful virtue, which is that it follows the way in which knowledge is actually produced and checked. Each discipline has its own standards of evidence, its own instruments, and its own history of error, and a proposal that crosses into a discipline should be subject to that discipline's standards and not to the standards of the proposer. This chapter surveys the principal disciplines and states, for each, the kinds of problem that are open, the form in which a precisely stated version of such a problem would appear, and the sense in which greater computational capability might bear on it. The survey is not exhaustive and makes no claim to settle the facts of any field. Its purpose is to illustrate the difference between a problem and the name of a problem, a difference on which the remainder of the book depends.

In physics and energy engineering the open problems are of two broad kinds. The first concerns the supply of energy at scale without emissions of carbon dioxide, and here the obstacles are largely engineering ones, since the relevant physics of photovoltaic conversion, wind capture, fission, and geothermal extraction is well established and the principal constraints are cost, materials, siting, grid integration, and the time required to build. The second concerns phenomena for which the underlying science is incomplete or for which the engineering is at the limit of what has been demonstrated, of which controlled thermonuclear fusion is the best known. Fusion is instructive because it is among the few problems whose requirements were stated with mathematical precision at an early date. Lawson's analysis of 1957~\cite{lawson1957} established that a deuterium--tritium plasma must satisfy a condition on the product of density, temperature, and energy confinement time for the fusion power released to exceed the power lost, and the condition has the form of an inequality that any proposed device either meets or fails to meet. The difficulty is not that the problem is vaguely specified but that satisfying the specification in a device that is also durable, maintainable, and economical remains unachieved. A capability that assists in the design of magnets, in the control of plasma instabilities, or in the discovery of materials that withstand neutron flux would be relevant to the problem, but it would be relevant by contributing to a precisely stated engineering target and not by pronouncing that the problem is solved. The comparison with grid-scale storage is similar: the target is a cost per unit of energy stored and delivered over a stated number of cycles, and the open questions concern chemistries, manufacturing, and degradation.

In chemistry and materials science the inventory includes catalysts and processes for the fixation of nitrogen at lower energy cost than the Haber--Bosch process, sorbents and processes for the capture of carbon dioxide from dilute sources, methods for the destruction or immobilization of persistent organic pollutants, and materials for batteries, membranes, and structures. A feature of this domain is that the space of candidate compounds is large beyond any possibility of exhaustive examination and that the relation between structure and property is determined by quantum mechanical calculations whose cost grows rapidly with system size. Machine learning methods that approximate these calculations, or that propose candidates for synthesis, can reduce the cost of searching the space, and they have been used for this purpose. The limit of the method is the one emphasized in Chapter 7, namely that a predicted compound must be synthesized and measured, and that the failure rate between prediction and confirmation determines the value of the prediction. In this domain the precise statement of a problem usually takes the form of a target property under stated conditions together with a cost ceiling, for example a catalyst that converts a stated feedstock at a stated temperature and pressure with a stated selectivity and lifetime.

In biology and medicine the inventory is dominated by problems whose difficulty is the complexity of the system and the ethical constraints on experimentation. Antimicrobial resistance threatens the efficacy of existing drugs and the economics of developing new ones. Many infectious diseases lack vaccines, and the design of vaccines for pathogens that mutate rapidly remains difficult. The prediction of the three-dimensional structure of proteins from sequence has advanced greatly, but the prediction of function, of interaction, and of the response of a whole organism to an intervention has not advanced in the same degree. The biology of aging, taken up in Chapter 27, belongs to this domain but is distinguished by the fact that its objectives are contested. In every case the decisive stages of validation are experimental and clinical, and are constrained by the time that biological processes require and by the protection owed to participants in trials. Computation shortens the preclinical stages of the pipeline; it does not shorten the clinical ones, and any claim that it does must be examined for hidden assumptions about regulatory and ethical processes.

In the earth and environmental sciences the open problems include the sensitivity of the climate to forcing at regional scales, the response of ecosystems and the carbon cycle to warming, the dynamics of ice sheets, the hydrology of aquifers and river basins, and the processes by which soils accumulate and lose carbon and fertility. These are problems of observation and modelling, and their difficulty derives from the number of interacting processes, the range of scales, and the limited length of the observational record. Improvements in computation permit models of higher resolution and larger ensembles, and methods that learn parameterizations from data promise to reduce errors in the representation of processes that cannot be resolved. The caution that attends this domain is that the verification of a model of a unique system such as the Earth cannot rely on repeated independent trials, so that confidence rests on out-of-sample performance on past records, on agreement among structurally different models, and on the physical consistency of the representation. Chapter 8 discusses the corresponding limits on prediction.

Mathematics and theoretical computer science form a distinct case because their problems are stated formally and their solutions can in principle be checked formally. A proof is either valid or not, and the development of proof assistants has made it possible to verify long arguments mechanically, so that the verification of a mathematical result may be nearly independent of the process by which it was found. This feature makes mathematics the domain in which the separation of generation from verification, defended in Chapter 6 and Chapter 10, is most completely realized, and it is a reason for optimism about the contribution of capable systems to the domain. It is also a reason for caution in extrapolating from it, because the problems of other disciplines lack a formal checker, and the success of mechanical verification in mathematics does not transfer to domains in which verification requires an experiment, a field trial, or a political settlement.

The social sciences and the study of institutions occupy a different position again. The open problems here include the design of mechanisms for allocating scarce goods, the design of organizations that remain accountable as they scale, the conditions under which commons are sustained, the dynamics of misinformation, and the measurement of well-being. The work of Ostrom~\cite{ostrom1990} on the governance of common-pool resources shows that durable arrangements exist and that they share features such as clearly defined boundaries, monitoring by participants, graduated sanctions, and accessible means of resolving conflicts. These are results about what works, obtained by comparative study, and they bear on every large technical program because every large technical program creates a common-pool problem. The relevance of computation is real in the form of simulation, data analysis, and the exploration of mechanism designs, but the legitimacy of an institution is conferred by the participation of those it governs, and no calculation substitutes for that participation.

Two general observations emerge from the survey. First, in every discipline the problems that most plausibly benefit from increased capability are those that have been stated precisely enough to admit a test, and the effort of stating a problem in that form is itself a considerable part of the work. Second, the disciplines differ in the cost and speed of verification, and this difference, more than any difference in the difficulty of generating candidate answers, determines how much of a given field can be accelerated. A field in which verification is cheap and fast, such as the checking of a formal proof, can absorb a large increase in the rate of generation. A field in which verification requires years of trial and the consent of participants cannot, and the proper response to this fact is to invest in verification capacity, in the form of laboratories, trial infrastructure, and independent institutions, and not to relax the standard of verification in order to match the new rate of generation.

\section*{Formalization}

The concept of a precisely stated problem used throughout the chapter can be made explicit. A \emph{well-posed problem} is a triple $P=(X,\varphi,c)$ where $X$ is a space of candidate solutions, $\varphi:X\to\{0,1\}$ is an acceptance predicate that can be evaluated for any candidate, and $c:X\to\R_{>0}$ is the cost of evaluating $\varphi$ at a candidate. A problem is \emph{open} if no known $x\in X$ has $\varphi(x)=1$. A \emph{search} for a solution is a sequence of candidates $x_1,x_2,\dots$, and its cost to first success is $\sum_{i\le\tau}c(x_i)$, where $\tau$ is the first index with $\varphi(x_\tau)=1$.

\begin{proposition}
Suppose candidates are generated at rate $g$ per unit time by a generator whose candidates satisfy $\varphi=1$ independently with probability $q>0$, and suppose each evaluation has cost $c$ and evaluation capacity $v$ evaluations per unit time. Then the expected time to the first accepted candidate is $\E[T]=1/(q\min(g,v))$. In particular, increasing $g$ above $v$ does not reduce $\E[T]$, and $\E[T]$ is inversely proportional to $q$ so that a tenfold increase in the quality of generation is equivalent to a tenfold increase in throughput when $g\le v$.
\end{proposition}

\begin{proof}
The number of candidates evaluated before the first success is geometric with mean $1/q$. Candidates can be evaluated at most at rate $\min(g,v)$, since none can be evaluated before it is generated and no more than $v$ can be evaluated per unit time. The expected time is the expected count divided by the rate.
\end{proof}

The proposition formalizes the second general observation of the chapter: acceleration of a field is limited by the slower of generation and verification, and the first lever to examine in any field is the evaluation rate $v$.

The fusion example gives a concrete acceptance predicate. In the deuterium--tritium case, Lawson's condition, in the form commonly quoted for ignition at temperatures near $10$--$20$ keV, requires the triple product $n\,T\,\tau_E\gtrsim3\times10^{21}\ \mathrm{keV\,s\,m^{-3}}$, where $n$ is the ion density, $T$ the ion temperature, and $\tau_E$ the energy confinement time. Here $X$ is the set of operating points together with devices that can sustain them, and the predicate is the stated inequality together with conditions of durability and net electrical output that the triple product alone does not capture. At fixed $T=15$ keV the required confinement time is $\tau_E\ge 3\times10^{21}/(nT)$, which gives $2$ s at $n=10^{20}\ \mathrm{m^{-3}}$, $4$ s at $n=5\times10^{19}\ \mathrm{m^{-3}}$, and $0.2$ s at $n=10^{21}\ \mathrm{m^{-3}}$. The elasticity of the required $\tau_E$ with respect to $n$ is $\partial\ln\tau_E/\partial\ln n=-1$, so that every tenfold increase of density permitted by the confinement scheme reduces the required confinement time tenfold, a trade-off that separates the magnetic and inertial approaches to the problem. The example shows what is meant by a problem having the form of an inequality that a device either satisfies or does not, and it also shows how much of the difficulty of the problem lies in the clauses not captured by the inequality, which is why the predicate $\varphi$ must be stated in full before any claim of progress can be assessed.
